\documentclass[10pt,a4paper]{amsart}
\usepackage{amsmath,amssymb,amsthm,graphicx}
\usepackage{verbatim}
\newenvironment{reference}{\comment}{\endcomment}
%\newenvironment{reference}{}{}
\newenvironment{slogan}{\comment}{\endcomment}
\newenvironment{history}{\comment}{\endcomment}

% For commutative diagrams we use Xy-pic
\usepackage[all]{xy}

% We use 2cell for 2-commutative diagrams.
\xyoption{2cell}
% Theorem environments.
%
\theoremstyle{plain}
\newtheorem{theorem}[subsection]{Theorem}
\newtheorem{proposition}[subsection]{Proposition}
\newtheorem{lemma}[subsection]{Lemma}

\theoremstyle{definition}
\newtheorem{definition}[subsection]{Definition}
\newtheorem{example}[subsection]{Example}
\newtheorem{exercise}[subsection]{Exercise}
\newtheorem{situation}[subsection]{Situation}

\theoremstyle{remark}
\newtheorem{remark}[subsection]{Remark}
\newtheorem{remarks}[subsection]{Remarks}

\numberwithin{equation}{subsection}

% Macros
%
\def\lim{\mathop{\mathrm{lim}}\nolimits}
\def\colim{\mathop{\mathrm{colim}}\nolimits}
\def\Spec{\mathop{\mathrm{Spec}}}
\def\Hom{\mathop{\mathrm{Hom}}\nolimits}
\def\Ext{\mathop{\mathrm{Ext}}\nolimits}
\def\SheafHom{\mathop{\mathcal{H}\!\mathit{om}}\nolimits}
\def\SheafExt{\mathop{\mathcal{E}\!\mathit{xt}}\nolimits}
\def\Sch{\mathit{Sch}}
\def\Mor{\mathop{\mathrm{Mor}}\nolimits}
\def\Ob{\mathop{\mathrm{Ob}}\nolimits}
\def\Sh{\mathop{\mathit{Sh}}\nolimits}
\def\NL{\mathop{N\!L}\nolimits}
\def\CH{\mathop{\mathrm{CH}}\nolimits}
\def\proetale{{pro\text{-}\acute{e}tale}}
\def\etale{{\acute{e}tale}}
\def\QCoh{\mathit{QCoh}}
\def\Ker{\mathop{\mathrm{Ker}}}
\def\Im{\mathop{\mathrm{Im}}}
\def\Coker{\mathop{\mathrm{Coker}}}
\def\Coim{\mathop{\mathrm{Coim}}}

% Boxtimes
%
\DeclareMathSymbol{\boxtimes}{\mathbin}{AMSa}{"02}

%
% Macros for moduli stacks/spaces
%
\def\QCohstack{\mathcal{QC}\!\mathit{oh}}
\def\Cohstack{\mathcal{C}\!\mathit{oh}}
\def\Spacesstack{\mathcal{S}\!\mathit{paces}}
\def\Quotfunctor{\mathrm{Quot}}
\def\Hilbfunctor{\mathrm{Hilb}}
\def\Curvesstack{\mathcal{C}\!\mathit{urves}}
\def\Polarizedstack{\mathcal{P}\!\mathit{olarized}}
\def\Complexesstack{\mathcal{C}\!\mathit{omplexes}}
% \Pic is the operator that assigns to X its picard group, usage \Pic(X)
% \Picardstack_{X/B} denotes the Picard stack of X over B
% \Picardfunctor_{X/B} denotes the Picard functor of X over B
\def\Pic{\mathop{\mathrm{Pic}}\nolimits}
\def\Picardstack{\mathcal{P}\!\mathit{ic}}
\def\Picardfunctor{\mathrm{Pic}}
\def\Deformationcategory{\mathcal{D}\!\mathit{ef}}
\usepackage{cleveref}
\crefname{originallemma}{original lemma}{original lemmas}
\Crefname{originallemma}{Original lemma}{Original lemmas}
\makeatletter
\newlabel{lemma-locally-ringed-intrinsic}{{04H7}{source}{}{sourcepointer-1}{}}
\newlabel{lemma-locally-ringed-intrinsic@cref}{{[originallemma][0][]04H7}{[1][0][]source}{}{}{}}
\makeatother
\title{Pullback of the unit group sheaf and intrinsicness of the locally ringed condition for morphisms}
\author{The Stacks Project Authors}
\date{stacks-project:sites-modules.tex:local-snapshot-20261012; GNU FDL 1.2 or later}
\begin{document}
\maketitle
\noindent\textit{Editorial note.} All delivered mathematical text is the author's own, copied verbatim from the pinned Stacks Project chapter file sites-modules.tex; no statement or proof step was rewritten, reordered or removed.\par\medskip
\noindent\textit{Editorial note.} The proof contains commutative diagrams typeset with Xy-pic; the retained two-slice preamble of this source supplies the native Xy-pic setup, so they compile as in the source.\par\medskip
\noindent\textit{Allowed standard prerequisites (editorial).}\begin{itemize}
\item Exactness of the pullback functor on sheaves. Stacks Project, Sites, section on morphisms of topoi (exactness of f\textasciicircum{}\{-1\}). Preserves the cartesianity of the diagram defining the units.
\item Definition of a morphism of locally ringed topoi. Stacks Project, Modules on Sites, section on locally ringed topoi (the defining diagram). Supplies the condition verified in the second half of the lemma.
\end{itemize}\medskip
\ifdefined\hypertarget\hypertarget{sourcepointer-1}{}\fi
\noindent\textit{Source reference.} Original lemma 04H7: sources/stacks-chow/tags.txt, source lines 5858--5858. Cross-chapter standard prerequisite cited by the author inside the counted span; it is used as a pointer and its own proof is not part of this item.\par\medskip
\begin{lemma}
\label{lemma-locally-ringed-intrinsic-morphism}
Let $f : \Sh(\mathcal{C}') \to \Sh(\mathcal{C})$
be a morphism of topoi. If $\mathcal{O}$ is a sheaf of rings
on $\mathcal{C}$, then
$$
f^{-1}(\mathcal{O}^*) = (f^{-1}\mathcal{O})^*.
$$
In particular, if $\mathcal{O}$ turns $\mathcal{C}$ into a locally
ringed site, then setting $f^\sharp = \text{id}$
the morphism of ringed topoi
$$
(f, f^\sharp) :
(\Sh(\mathcal{C}'), f^{-1}\mathcal{O})
\to
(\Sh(\mathcal{C}, \mathcal{O})
$$
is a morphism of locally ringed topoi.
\end{lemma}

\begin{proof}
Note that the diagram
$$
\xymatrix{
\mathcal{O}^* \ar[rr] \ar[d]_{u \mapsto (u, u^{-1})} & &
{*} \ar[d]^{1} \\
\mathcal{O} \times \mathcal{O} \ar[rr]^-{(a, b) \mapsto ab} & &
\mathcal{O}
}
$$
is cartesian. Since $f^{-1}$ is exact we conclude that
$$
\xymatrix{
f^{-1}(\mathcal{O}^*)
\ar[d]_{u \mapsto (u, u^{-1})} \ar[rr] & &
{*} \ar[d]^{1} \\
f^{-1}\mathcal{O} \times f^{-1}\mathcal{O} \ar[rr]^-{(a, b) \mapsto ab} & &
f^{-1}\mathcal{O}
}
$$
is cartesian which implies the first assertion. For the second,
note that $(\mathcal{C}', f^{-1}\mathcal{O})$ is a locally ringed site
by
Lemma \ref{lemma-locally-ringed-intrinsic}
so that the assertion makes sense. Now the first part implies that
the morphism is a morphism of locally ringed topoi.
\end{proof}

\end{document}
